\uuid{h2bF}
\exo7id{1695}
\titre{exo7 1695}
\auteur{bodin}
\organisation{exo7}
\datecreate{1999-11-01}
\isIndication{false}
\isCorrection{false}
\chapitre{Réduction d'endomorphisme, polynôme annulateur}
\sousChapitre{Réduction de Jordan}
\cpge{mpsi-mp2i.polynomes-fractions-rationnelles.b, pcsi.polynomes.b, ptsi.polynomes.b, mpsi-mp2i.applications-lineaires.a, pcsi.applications-lineaires.a, ptsi.applications-lineaires.a, mpsi-mp2i.applications-lineaires.b, pcsi.applications-lineaires.b, ptsi.applications-lineaires.b, bcpst1.applications-lineaires-matrices, mp-mpi.reduction-endomorphismes.d, pc.reduction.c, psi.reduction-endomorphismes.c, pt.reduction-endomorphismes.c, mp-mpi.reduction-endomorphismes.i, pc.reduction.d, psi.reduction-endomorphismes.d, bcpst2.valeurs-propres-vecteurs-propres.a, bcpst2.valeurs-propres-vecteurs-propres.b}
\module{Algèbre}
\niveau{L2}
\difficulte{}

\contenu{
\texte{
Soit $ \rho  $ l'application de $ \R_4[X] $ dans
lui-m\^eme qui \`a un polyn\^ome $ P  $ associe le reste de la
division euclidienne de $ P $ par $ (X^2-1) .$
}
\begin{enumerate}
    \item \question{Montrer que $ \rho  $ est lin\'eaire.}
    \item \question{Montrer que $ \rho ^2= \rho .$ En d\'eduire que $
\rho  $ est diagonalisable.}
    \item \question{D\'eterminer (de pr\'ef\'erence
sans calcul) une base de vecteurs propres pour $ \rho  .$}
\end{enumerate}
}
